\chapter{Chặng đường phía trước}
\label{ch:roadahead}

Với frontend được phục vụ từ cụm (Chương~26), mọi thứ có thể triển
khai của dự án đều đã được triển khai và ứng dụng dùng được trong
trình duyệt. Các cột mốc còn lại, mỗi cái thay một thứ tự chế bằng cơ
chế công nghiệp --- và mỗi cái đều có ý nghĩa nhất vào lúc này, khi
bạn đã biết nó thay thế cái gì.

\section{Khả năng quan sát}

Bên ngoài \code{kubectl logs} là bộ ba tiêu chuẩn: \textbf{metric}
(Prometheus cào các endpoint \code{/actuator/prometheus} mà mọi
service đã mở sẵn; dashboard Grafana trên bộ ba RED --- rate, errors,
duration), \textbf{trace} (dây nối Zipkin đã có trong mọi cấu hình,
hiện lấy mẫu ở mức không trong cụm --- một biến môi trường bật nó lên
khi có pod Zipkin), và \textbf{alert} (Alertmanager: một cú báo khi tỷ
lệ lỗi của một service hay độ đầy của một PVC vượt ngưỡng). Việc này
đến qua Helm chart \code{kube-prometheus-stack} --- khoảnh khắc đã hứa
khi bạn \emph{dùng} Helm mà không tự viết chart (hộp quyết định của
Chương~20, được giữ lời). Strimzi có thể xuất metric của broker vào
cùng Prometheus đó với vài dòng \code{metricsConfig} trên resource
Kafka --- độ trễ của consumer group là alert Kafka đầu tiên đáng có.

\begin{hardware}
Bộ đầy đủ tốn $\sim$1,5\,GB và không ít CPU --- ngân sách nói là vừa,
nhưng đây là thứ đầu tiên phải cắt khi tải cao. Khả năng quan sát có
thể bỏ đi là món xa xỉ của homelab; các đội production tính kích cỡ
cho nó trước tiên.
\end{hardware}

\section{GitOps: Argo CD}

Pipeline \emph{đẩy} các lần triển khai (Jenkins giữ một kubeconfig và
chạy kubectl). GitOps đảo chiều mũi tên: một agent trong cụm (Argo CD)
\emph{kéo} --- liên tục so sánh manifest của một kho git với cụm và
hội tụ. Điều đó mua được gì, bằng vốn từ bạn đã có: thông tin xác
thực của cụm không bao giờ rời khỏi cụm (nỗi lo RBAC của Chương~22
tan biến --- CI chỉ push image và commit một lần nâng tag); việc xóa
cuối cùng cũng lan truyền (cạm bẫy prune của Chương~20, được giải
quyết từ thiết kế); trôi (drift) cấu hình trở nên hữu hình (một
Deployment bị sửa tay hiện ``OutOfSync''); và rollback là \code{git
revert}. Jenkins giữ CI --- build, test, push, cập nhật tag trong git
--- Argo CD nhận CD. Stage deploy của Jenkinsfile, tức Chương~23 của
cuốn sách này, trở thành một commit.

\section{Cho stack cổ điển về hưu}

Sự dư thừa có chủ đích của Chương~2 và 3 nay có thể được thanh toán,
từng service một:

\begin{itemize}
  \item \textbf{Bỏ Eureka}: đổi các route của gateway từ
  \code{lb://course-service} sang \code{http://course-service:8082} ---
  DNS của Service và các endpoint gác bằng readiness làm mọi thứ
  registry từng làm ở đây. Xóa một thứ triển khai, giải phóng
  512\,Mi.
  \item \textbf{Bỏ Config-server}: mount các file
  \code{configurations/} thành một ConfigMap và trỏ
  \code{spring.config.import} vào thư mục đã mount
  (\code{optional:file:/config/}). Cùng file, bớt một service, và thay
  đổi cấu hình trở thành thay đổi manifest chảy qua GitOps.
\end{itemize}

Cả hai cuộc di dời đều là điểm cộng cho CV chính vì chúng chạy
\emph{về phía} nền tảng: biết vì sao Eureka tồn tại \emph{và} vì sao
Kubernetes khiến nó lỗi thời là khác biệt giữa dùng một trong hai và
hiểu cả hai.

\section{Jenkins vào cụm}

Phép đảo ngược kiểu homelab của Chương~22 --- build trên controller,
Jenkins nằm ngoài k3s --- đảo lại gọn gàng một khi pipeline đã ổn
định: plugin Kubernetes sinh một pod agent cho mỗi build, và bản thân
Jenkins trở thành một StatefulSet với \code{JENKINS\_HOME} trên một
PVC. Cuộc di dời chạm tới mọi khái niệm của Phần~V cùng lúc, đó là lý
do nó được xếp cuối cùng.

\section{OpenShift, nói ngắn}

Các manifest được giữ khả chuyển có chủ đích: image không chạy bằng
root của Jib (Chương~13) thỏa mãn mô hình bảo mật UID tùy ý của
OpenShift, nên triển khai lên một sandbox OpenShift miễn phí hẳn chỉ
cần bước dịch Ingress sang Route (Chương~18). Chạy thí nghiệm đó một
lần sẽ kiểm chứng những tuyên bố về tính khả chuyển mà cuốn sách này
vẫn đưa ra --- và vốn từ của OpenShift (Route, SCC, BuildConfig) ánh
xạ một-một lên các chương bạn đã đọc.

\section{Bạn đang đứng ở đâu}

Bỏ tên dự án đi, bản kiểm kê đọc như một bản mô tả công việc
platform: bộ máy phân giải thuộc tính giúp một artifact chạy ở bất cứ
đâu; log, consumer group và ngữ nghĩa giao nhận; image, layer và
registry; điều hòa, probe, RBAC, storage class và operator; pipeline,
danh tính deploy giới hạn bằng RBAC, và phép đảo chiều GitOps. Những
ràng buộc của homelab không phải là giới hạn --- chúng là giáo trình:
mỗi hộp ``trên phần cứng này'' đều ép ra một quyết định mà tiền bạc
thường che đi.

\begin{quiz}
\begin{enumerate}
  \item Bạn sẽ đặt alert trên metric nào trước tiên cho đường Kafka,
  và nó đo cái gì theo cách nói của Chương~10?
  \item Triển khai kiểu đẩy so với kiểu kéo: chỉ ra thông tin xác thực
  của cụm nằm ở đâu trong mỗi kiểu, và giải thích việc dọn (prune) của
  GitOps giải quyết cạm bẫy của chương nào.
  \item Phác thảo việc cho Eureka về hưu: cái gì đổi trong cấu hình
  của gateway, hai cơ chế Kubernetes nào thay cho đăng ký và liveness,
  và cái gì có thể xóa?
  \item Vì sao image của Jib khiến thí nghiệm OpenShift rẻ, và phép
  dịch object nào vẫn còn lại?
\end{enumerate}
\end{quiz}
